%! Author = hilfiker
%! Date = 21.07.2022

% Preamble
\documentclass[11pt]{article}

% Packages
\usepackage{amsmath} % Formulas in document
\usepackage{pdflscape} % Landscaped pages in Pdf
\usepackage[a4paper, margin=0.8in]{geometry} % Set the margin and size of a page
\usepackage[hidelinks]{hyperref} % Remove Boxes around Hyperlinks
\usepackage{lastpage} % Custom page numbering
\usepackage{fancyhdr}
\usepackage{graphicx}
\usepackage{tcolorbox} % Custom page numbering

% Configure page numbering & footer
\pagestyle{fancy}
\fancyhf{}
\lfoot{Tobias Hilfiker}
\cfoot{Seite \thepage \hspace{1pt} von \pageref{LastPage}}
\rfoot{\today}

% Methadata
\title{
    \includegraphics{media/frontpage.jpg}
    Modul 152 \\
    Multimediainhalte in einen Webauftritt integrieren}
\author{Tobias Hilfiker}
\date{\today}

% Document
\begin{document}

    %Kapitel Frontpage
    \maketitle
    \pagebreak

    %Kapitel Table of contents
    \tableofcontents
    \pagebreak

    %Kapitel 1 - Einleitung
    \section{Einleitung}

    %Kapitel 2 - Thema
    \section{Thema}

    todo kleiner Text

    %Kapitel 2.1 - Wahl des Themas
    \subsection{Wahl des Themas}

    Ich habe das Thema Curling ausgewählt, da es bereits seit vielen Jahren mein Hobby ist. Jeden Winter
    von ca. September bis April spiele ich in einem Juniorenteam der Liga C in St.Gallen. Curling eignet
    sich für dieses Modul sehr, da in der Halle immer die gleichen Lichtverhältnisse herrschen. Da in
    der Halle in St. Gallen auch internationale Tourniere ausgetragen werden, ist die Beleuchtung bereits
    sehr gut.

    %Kapitel 3 - Visualisierung E-Portfolio
    \section{Visualisierung E-Portfolio}

    %Kapitel 3.1 - Menüstruktur
    \subsection{Menüstruktur}
    Als Menü-Navigation habe ich mich für eine Navigation am oberen Teil der Webseite

    %Kapitel 3.2 - Farbwahl
    \subsection{Farbwahl}

    %Kapitel 3.3 - Mockups
    \subsection{Mockups}
    - Mockups mit Adobe xd -> Wieso
    - Bilder der Mockups einfügen

    %Kapitel 4 - Bildbearbeitung
    \section{Bildbearbeitung}

    %Kapitel 4.1 - Bildbearbeitungsmethode
    \subsection{Bildbearbeitungsmethode}

    %Kapitel 5 - Logo
    \section{Logo}

    %Kapitel 5.1 - Idee
    \subsection{Idee}

    %Kapitel 6 - Technisches Konzept E-Portfolio
    \section{Technisches Konzept E-Portfolio}
    -Welche Technologien werden verwendet
    -

    %Kapitel 7 - Reflexion
    \section{Reflexion}

    %Kapitel 7.1 - Bildbearbeitung
    \subsection{Bildbearbeitung}

    %Kapitel 7.2 - Logodesign
    \section{Logodesign}


\end{document}